\documentclass{article}
\usepackage[margin=28mm]{geometry}
\usepackage{amsmath,amssymb}
\usepackage{xurl}
\usepackage[hidelinks]{hyperref}
% Shared commands for notes in docs/paper-gaps/.

\DeclareUnicodeCharacter{2080}{\ifmmode{}_{0}\else\textsubscript{0}\fi}
\DeclareUnicodeCharacter{2081}{\ifmmode{}_{1}\else\textsubscript{1}\fi}
\DeclareUnicodeCharacter{2082}{\ifmmode{}_{2}\else\textsubscript{2}\fi}
\DeclareUnicodeCharacter{2099}{\ifmmode{}_{n}\else\textsubscript{n}\fi}

\newcommand{\Lean}{\textsc{Lean}}
\ExplSyntaxOn
\NewDocumentCommand{\leanid}{m}{
  \ifmmode
    \textsf{\detokenize{#1}}
  \else
    \group_begin:
    \urlstyle{sf}
    \tl_set:Nn \l_tmpa_tl {#1}
    \tl_replace_all:Nnn \l_tmpa_tl {\_} {_}
    \exp_args:NV \path \l_tmpa_tl
    \group_end:
  \fi
}
\ExplSyntaxOff
\providecommand{\tr}{\operatorname{tr}}
\newcommand{\tnleanIssueBase}{https://github.com/LionSR/TNLean/issues/}
\newcommand{\tnleanPrBase}{https://github.com/LionSR/TNLean/pull/}
\newcommand{\tnleanDocBase}{https://sirui-lu.com/TNLean/docs/}
\newcommand{\ghissue}[1]{\href{\tnleanIssueBase#1}{GitHub issue \##1}}
\newcommand{\ghpr}[1]{\href{\tnleanPrBase#1}{PR \##1}}
\newcommand{\doclink}[1]{\href{\tnleanDocBase#1.html}{\path{#1}}}

% Dirac notation and the identity operator, as in blueprint/src/macros/common.tex.
% \providecommand keeps note-local definitions authoritative.
\usepackage{dsfont}
\providecommand{\ket}[1]{|#1\rangle}
\providecommand{\bra}[1]{\langle #1|}
\providecommand{\braket}[2]{\langle #1 | #2 \rangle}
\providecommand{\ketbra}[2]{|#1\rangle\!\langle #2|}
\providecommand{\Id}{\mathds{1}}
\providecommand{\one}{\mathds{1}}
\providecommand{\C}{\mathbb{C}}
\providecommand{\MN}[1]{M_{#1}(\C)}
\providecommand{\MD}{M_D(\C)}

% Verdict marker: \gapnote{<kind>}{<status>}. Kind is the policy
% classification (clarification, local-correction, scope-restriction,
% unfaithful, false-source, open-gap); status is open, wip (elimination
% underway), resolved, or historical. Machine-read by the site index;
% prints a verdict line.
\newcommand{\gapnote}[2]{\par\noindent\textbf{Verdict:} #1 (#2).\par}

\providecommand{\gapnote}[2]{\par\noindent\textbf{Verdict:} #1 (#2).\par}
\title{The Final State in Hastings' Energy Estimate (hchy)}
\author{}
\date{}
\begin{document}
\maketitle
\gapnote{local-correction}{resolved}

\section{The asserted estimate}
Hastings, \emph{Solving gapped Hamiltonians locally}~\cite{Hastings2006Solving},
\href{https://arxiv.org/abs/cond-mat/0508554v3}{arXiv:cond-mat/0508554v3},
source \path{proj3new.tex}, lines 705--720, states that after the last of the
$h=\log_2N$ levels of the hierarchical construction ``there is at least one
state $\alpha^h$, and one may show''
\begin{align}
 \langle\alpha^h,\textstyle\sum_iM_i\,\alpha^h\rangle
 \le\langle\Psi_a,\textstyle\sum_iM_i\,\Psi_a\rangle
 +NJ\,O\bigl(e^{-l_{\rm proj}/(\log_2(N)\,l_{\rm fac})}\bigr).
 \label{eq:hastings05_hierarchical_energy_hchy}
\end{align}
No argument is given.

\section{What is proved}
Let $Q$ be the projection onto the states kept at the last level and suppose
$\|\Psi_a-Q\Psi_a\|\le b<1$, the norm form of the deficiency used for every
level. Then $\phi=Q\Psi_a/\|Q\Psi_a\|$ is a unit vector with $Q\phi=\phi$ and
$\|\phi-\Psi_a\|\le2b$, and for every operator $M$
\begin{align}
 \operatorname{Re}\langle\phi,M\phi\rangle
 \le\operatorname{Re}\langle\Psi_a,M\Psi_a\rangle+2\|M\|\,\|\phi-\Psi_a\|
 \le\operatorname{Re}\langle\Psi_a,M\Psi_a\rangle+4\|M\|b,
 \label{eq:hastings05_hierarchical_energy_proved}
\end{align}
by writing the difference of the two expectations as
$\langle\phi-\Psi_a,M\phi\rangle+\langle\Psi_a,M(\phi-\Psi_a)\rangle$. For
$M=\sum_{i<N}M_i$ with $\|M_i\|\le J$ the error term is $4NJb$, so with $b=b_h$
(\ref{eq:hastings05_hierarchical_energy_proved}) is
(\ref{eq:hastings05_hierarchical_energy_hchy}) with the explicit constant
$4NJb_h$.\footnote{\leanid{TNLean.MPS.AreaLaw.LocalSolution.exists_mem_range_re_inner_le_add}
    and
    \leanid{TNLean.MPS.AreaLaw.LocalSolution.exists_mem_range_re_inner_sum_le_add}
    in \doclink{TNLean/MPS/AreaLaw/LocalSolution/HierarchicalEnergy}.}

\section{Every final state}
At level $m+1$ the source compresses the operator $M$ joining the two halves
of a block to the range of the projection $Q$ of the level below and keeps
the eigenvectors of $QMQ$ with eigenvalue within $x_m$ of
$c=\langle\Psi_a,M\Psi_a\rangle$ (lines 640--642 and 671--685). A unit vector
$\phi$ kept by the final projection is kept at every level. For each joining
operator it satisfies $Q\phi=\phi$ and lies in the window of $QMQ$, where
$\|(QMQ-c)\phi\|\le x_m$, so
\begin{align}
 |\operatorname{Re}\langle\phi,M\phi\rangle-c|
 =|\operatorname{Re}\langle\phi,(QMQ-c)\phi\rangle|\le x_m.
 \label{eq:hastings05_hierarchical_energy_window}
\end{align}
Summing over the $2^{h-1-m}$ joining operators of each level $m$ gives
\begin{align}
 \operatorname{Re}\langle\phi,M_{\rm open}\,\phi\rangle
 \le\operatorname{Re}\langle\Psi_a,M_{\rm open}\Psi_a\rangle
 +\sum_{m<h}2^{h-1-m}x_m,
 \label{eq:hastings05_hierarchical_energy_every_state}
\end{align}
where $M_{\rm open}$ is the sum of the joining operators. This holds for every
unit vector in the final space, in particular for each
$\alpha^h$.\footnote{\leanid{TNLean.MPS.AreaLaw.LocalSolution.abs_re_inner_hierarchySum_sub_le}
    and
    \leanid{TNLean.MPS.AreaLaw.LocalSolution.re_inner_hierarchySum_le_add}
    in \doclink{TNLean/MPS/AreaLaw/LocalSolution/FinalStateEnergy}.}
Cross terms between different $\alpha^h$ play no role, and together with the
count of low-energy orthonormal states it gives at most $2n$ final states
once the right side of
(\ref{eq:hastings05_hierarchical_energy_every_state}) is within $\Delta E/2$
of the ground energy.\footnote{\leanid{TNLean.MPS.AreaLaw.LocalSolution.card_le_two_mul_re_trace_of_hierarchy}
    in \doclink{TNLean/MPS/AreaLaw/LocalSolution/FinalStateEnergy}.}

\section{The closing link}
The chain of supersites is periodic (lines 626--628). The binary tree of joining
operators covers the $N-1$ links $(k,k+1)$, $1\le k<N$, of the open chain,
and the link $(N,1)$ joins the two halves of no block, so no window of the $h$
levels controls its energy. Bound~(\ref{eq:hastings05_hierarchical_energy_every_state})
therefore concerns $M_{\rm open}=\sum_iM_i-M_{N,1}$, not $\sum_iM_i$, and a count
of final states drawn from it would need the gap of $M_{\rm open}$, whereas the
source assumes a gap for the original Hamiltonian $H$. Neither the open
sum nor the truncated periodic sum is given a separate gap. Bounding the
missing term by its norm adds
about $2\|M_{N,1}\|$, of order $Jl_{\rm proj}$, which is far larger than the
error in (\ref{eq:hastings05_hierarchical_energy_hchy}) and than the budget
$\Delta E/2$ of the count. The source does not address this term.

The fix is one further window. Let $Q=P^h_0$ be the final projection of the
tree, and let $R$ project onto the eigenvectors of $QM_{N,1}Q$ in the range of
$Q$ with eigenvalue within $x_h$ of $\langle\Psi_a,M_{N,1}\Psi_a\rangle$. This
is the step of a further level whose second half is the whole space, so the
one-level estimate gives
\begin{align}
 \|\Psi_a-R\Psi_a\|^2\le2D_h^2+\bigl((\varepsilon+2KD_h)/x_h\bigr)^2=D_{h+1}^2,
 \label{eq:hastings05_hierarchical_energy_closing_deficiency}
\end{align}
with $D_m$ the deficiency recursion of the tree, $\varepsilon$ the residual
and $K$ the norm bound of the links. With the source's radii,
$x_m=K(e+2D_m)^{1-s}$ for $\varepsilon=Ke$ and $s=1-1/h$, the closed form of the
deficiencies applies at every level, the extra one included:
$D_{h+1}\le(3^{h+1}-1)e^{s^{h+1}}$. Since $s^{h+1}=(1-1/h)\,s^h$ and
$s^h\ge1/4$ for $h\ge2$, this is at most $3\cdot3^h e^{1/8}$ for $e\le1$, of the
same kind as the bound $(3^h-1)e^{1/4}$ of the last level of the tree: polynomial in
$N=2^h$ times a power of the residual $e$, which is exponentially small in
$l_{\rm proj}/l_{\rm fac}$. Under the source's choice of $l_{\rm proj}$ it is
therefore below $1$, so $R\Psi_a\ne0$ and at least one final state
survives.\footnote{\leanid{TNLean.MPS.AreaLaw.LocalSolution.norm_sub_closingLinkWindow_le}
    in \doclink{TNLean/MPS/AreaLaw/LocalSolution/ClosingLink}.}
A unit vector $\phi$ with $R\phi=\phi$ has $Q\phi=\phi$ and
$|\operatorname{Re}\langle\phi,M_{N,1}\phi\rangle-\langle\Psi_a,M_{N,1}\Psi_a\rangle|\le x_h$
by (\ref{eq:hastings05_hierarchical_energy_window}), so adding this to
(\ref{eq:hastings05_hierarchical_energy_every_state}) gives
\begin{align}
 \operatorname{Re}\langle\phi,\textstyle\sum_iM_i\,\phi\rangle
 \le\operatorname{Re}\langle\Psi_a,\textstyle\sum_iM_i\,\Psi_a\rangle
 +\sum_{m<h}2^{h-1-m}x_m+x_h,
 \label{eq:hastings05_hierarchical_energy_periodic}
\end{align}
which is eq.~(\ref{eq:hastings05_hierarchical_energy_hchy}) for the periodic
sum. If the periodic sum has a gap, this gives its abstract state-count
bound; the physical application instead compares that sum with the original
Hamiltonian $H$, as below.\footnote{\leanid{TNLean.MPS.AreaLaw.LocalSolution.re_inner_hierarchySum_add_le_add}
    and
    \leanid{TNLean.MPS.AreaLaw.LocalSolution.card_le_two_mul_re_trace_of_closingLink}
    in \doclink{TNLean/MPS/AreaLaw/LocalSolution/ClosingLink}.}
The final states are thus those kept by $R$, a subspace of the source's
final space: the construction has one window more than the source states,
and its error has one radius more. Bound~(\ref{eq:hastings05_hierarchical_energy_periodic})
and the count of at most $2n$ states are proved for this corrected
construction, not for the source's own states $\alpha^h$, whose energy for the
link $(N,1)$ remains uncontrolled. The eigenvectors of $QM_{N,1}Q$ are linear
combinations of the $\alpha^h$ at the root. Every vector retained by $R$ lies
in the range of $Q$, hence in the retained spaces of both children, and
recursively in every retained space below the root. The closing window
therefore preserves the tree of nested local spectral projections. This is
the generalized matrix product representation asserted in the source
(lines 596--607 and 732--743); the source gives no bound on the dimensions
of the intermediate retained spaces.\footnote{\leanid{TNLean.MPS.AreaLaw.LocalSolution.IsRetainedHierarchy},
\leanid{TNLean.MPS.AreaLaw.LocalSolution.isRetainedHierarchy_of_hierarchyProjection}, and
\leanid{TNLean.MPS.AreaLaw.LocalSolution.isRetainedHierarchy_of_hierarchyClosingProjection}.}
The open-chain statements remain as the case without the closing link.
\section{Explicit construction and the physical gap}
Number the supersites from zero and let $a_j$ join supersites $j-1$ and $j$,
with $a_0$ the closing edge. The operator joining the halves of block $k$
at level $m+1$ is $a_{(2k+1)2^m}$. Induction by splitting the interval at its
midpoint gives
\begin{align}
 T_k^m(M)=\sum_{k2^m<j<(k+1)2^m}a_j,
 \qquad T_0^h(M)+a_0=\sum_{0\le j<2^h}a_j.
 \label{eq:hastings05_tree_edges}
\end{align}
The projections are defined recursively from these operators and their
compressed spectral windows; neither the projections nor their recursion
are additional assumptions. Their support follows from that of the
edges.\footnote{\leanid{TNLean.MPS.AreaLaw.LocalSolution.hierarchyProjection},
\leanid{TNLean.MPS.AreaLaw.LocalSolution.hierarchyClosingProjection},
\leanid{TNLean.MPS.AreaLaw.LocalSolution.hierarchySum_hierarchyEdge}, and
\leanid{TNLean.MPS.AreaLaw.LocalSolution.hierarchyEdge_mem_supportedOperators}.}

Write $T=\sum_j a_j$ and suppose $\|H-T\|\le\delta$, where $H$ is the
physical Hamiltonian with ground energy $E_0$ and gap $\Delta>0$ above its
ground projection $P_0$. For a unit ground vector $\psi$ and a unit vector
$\phi$ retained by the corrected final projection, the window bound gives
$\operatorname{Re}\langle\phi,T\phi\rangle
\le\operatorname{Re}\langle\psi,T\psi\rangle+\eta$, where
$\eta=T_0^h(x)+x_h$. Two uses of the operator-norm estimate give
\begin{align}
 \operatorname{Re}\langle\phi,H\phi\rangle
 &\le\operatorname{Re}\langle\psi,H\psi\rangle+2\delta+\eta
 =E_0+2\delta+\eta.
 \label{eq:hastings05_physical_energy}
\end{align}
Consequently $2\delta+\eta\le\Delta/2$ suffices for the count of at most
$2\operatorname{tr}P_0$ orthogonal retained states. No gap of $T$ is
assumed.\footnote{\leanid{TNLean.MPS.AreaLaw.LocalSolution.card_le_two_mul_re_trace_of_approximation}.}

\section{The pair bounds on an arbitrary periodic physical chain}

The centred bound needed by the quantitative hierarchy is derived from
individual local terms. Let the physical chain have $M\ge1$ sites and let
$\ell\ge1$ be the integer projection radius. Put
$N=\max\{1,\lfloor M/(2\ell)\rfloor\}$. If $M\ge2\ell$, take $N-1$
blocks of width $2\ell$ and include the remainder in the last block. Its
width is between $2\ell$ and $4\ell$. A shorter chain is a single block.
No physical site or Hilbert-space factor is added.

For a physical site with coordinate $z$, put
\begin{align}
 b(z)&=\min\{\lfloor z/(2\ell)\rfloor,N-1\},\qquad
 j(z)=\begin{cases}
 \lfloor(z+\ell)/(2\ell)\rfloor,&\lfloor(z+\ell)/(2\ell)\rfloor<N,\\
 0,&\text{otherwise}.
 \end{cases}
 \label{eq:hastings05_rounded_assignment}
\end{align}
A term $M_i$ supported in the physical cyclic ball of radius $\ell$ acts
on the two blocks adjacent to $j(i)$. This includes the edge joining the
last and first blocks. At most $4\ell$ centres are assigned to an edge.
Writing $a_j=\sum_{i:j(i)=j}M_i$, the total sum is unchanged. Thus
$\|M_i\|\le J$ and the source's local estimate
$\|(M_i-\langle M_i\rangle_\psi)\psi\|\le Je^{-r}$ give
\begin{align}
 \|a_j-\langle a_j\rangle_\psi\|&\le8\ell J,\\
 \|(a_j-\langle a_j\rangle_\psi)\psi\|&\le4\ell J e^{-r}.
 \label{eq:hastings05_actual_pair_bounds}
\end{align}
Both bounds are independent of $M$. For equal blocks, exactly $2\ell$
centres are assigned to an edge, improving the constants to $4\ell J$
and $2\ell Je^{-r}$.\footnote{Declarations:
\leanid{TNLean.MPS.AreaLaw.LocalSolution.roundedPeriodicBlockCount},
\leanid{TNLean.MPS.AreaLaw.LocalSolution.roundedPeriodicEdge_support_of_cover},
\leanid{TNLean.MPS.AreaLaw.LocalSolution.roundedPeriodicEdge_zero_support},
\leanid{TNLean.MPS.AreaLaw.LocalSolution.card_roundedPeriodicEdge_fibre_le},
\leanid{TNLean.MPS.AreaLaw.LocalSolution.sum_roundedPeriodicOperator},
\leanid{TNLean.MPS.AreaLaw.LocalSolution.residual_roundedPeriodicOperator_le}, and
\leanid{TNLean.MPS.AreaLaw.LocalSolution.centeredNorm_roundedPeriodicOperator_le}.}

\section{The physical approximation is a sum of filtered-term errors}

Let $H=\sum_iH_i$, and let $\widetilde H_i$ be the full normalized Gaussian
filter of $H_i$ under $H$. Filtering fixes $H$ and is additive, so
\begin{align}
 \sum_i\widetilde H_i=H.
 \label{eq:hastings05_full_filtered_sum}
\end{align}
The small local difference is $\widetilde H_i-M_i$, not $H_i-M_i$.
In particular the preceding localization estimate
$\|\widetilde H_i-M_i\|\le Je^{-r}$ implies
\begin{align}
 \Bigl\|H-\sum_j a_j\Bigr\|
 =\Bigl\|\sum_i(\widetilde H_i-M_i)\Bigr\|
 \le MJe^{-r}\le4\ell NJe^{-r}.
 \label{eq:hastings05_actual_total_error}
\end{align}
Here $J$ absorbs the fixed factors in the local norm, residual and
localization estimates. A Gaussian-filtered term need not be close in
operator norm to the corresponding unfiltered interaction.\footnote{
\leanid{TNLean.Lattice.sum_gaussianFilteredTerm} and
\leanid{TNLean.MPS.AreaLaw.LocalSolution.card_le_two_mul_re_trace_of_roundedPeriodicProjection}.}

\section{An explicit integer quadratic logarithmic radius}

Let $r=\eta\ell$ for a decay rate $\eta>0$. The actual centred bound in
(\ref{eq:hastings05_actual_pair_bounds}) is
$K=8\ell J=J_{\rm eff}r$, with $J_{\rm eff}=8J/\eta$, independently of the
physical chain length. Define
\begin{align}
 A_M&=512(\lceil\log_2M\rceil+2)^2
       \bigl(1+\max\{0,\log(J_{\rm eff}/\Delta)\}\bigr),\\
 \ell&=\max\{\lceil R\rceil,\lceil A_M/\eta\rceil\},\qquad r=\eta\ell.
 \label{eq:hastings05_integer_radius}
\end{align}
This integer is positive, is at least the physical interaction range, and
satisfies $\ell<\lceil R\rceil+A_M/\eta+1$.
For the block count determined by this same radius, put
$h=\max\{2,\lceil\log_2N\rceil\}$. Since $N\le M$, the quantitative
hierarchy length is at most $A_M\le r$. Thus the choice is explicit;
neither divisibility of $M$ by $2\ell$ nor an implicit equation for the
number of blocks is required.\footnote{Declarations:
\leanid{TNLean.MPS.AreaLaw.LocalSolution.roundedPeriodicProjectionRadius},
\leanid{TNLean.MPS.AreaLaw.LocalSolution.roundedPeriodicProjectionRadius_pos},
\leanid{TNLean.MPS.AreaLaw.LocalSolution.hierarchyProjectionLength_le_chosen_roundedPeriodicRadius}, and
\leanid{TNLean.MPS.AreaLaw.LocalSolution.roundedPeriodicProjectionRadius_lt},
\leanid{TNLean.MPS.AreaLaw.LocalSolution.physicalPeriodicProjectionRadius_spec},
\leanid{TNLean.MPS.AreaLaw.LocalSolution.hierarchyProjectionLength_le_physicalPeriodicProjectionRadius}, and
\leanid{TNLean.MPS.AreaLaw.LocalSolution.physicalPeriodicProjectionRadius_lt}.}

At this length the sum of all tree and closing-window radii is at most
$\Delta/4$. The scalar estimate also includes the linear prefactor:
\begin{align}
 2^hJ_{\rm eff}r e^{-r}\le\Delta/8.
 \label{eq:hastings05_linear_truncation_budget}
\end{align}
Together with (\ref{eq:hastings05_actual_total_error}), this gives
$\delta\le\Delta/8$, hence $2\delta+\sum_{m<h}2^{h-1-m}x_m+x_h\le\Delta/2$.
The local residuals construct a retained unit vector within
$2(3^{h+1}-1)e^{-r/8}$ of $\psi$, and the original physical gap gives
at most $2\dim\ker(H-E_0)$ orthonormal retained states.
Every retained vector remains in the space of every node of the tree.
The complete periodic construction therefore has range $O(\log^2M)$
at fixed local constants, decay rate and gap, with the corrected closing
window.\footnote{Declarations:
\leanid{TNLean.MPS.AreaLaw.LocalSolution.hierarchyTruncationError_mul_length_le_eighth_gap},
\leanid{TNLean.MPS.AreaLaw.LocalSolution.exists_unit_mem_roundedPeriodicProjection_of_chosen_radius}, and
\leanid{TNLean.MPS.AreaLaw.LocalSolution.card_le_two_mul_re_trace_of_roundedPeriodicProjection_of_chosen_radius}.}

The physical construction derives the localization estimate from the
original interactions. The cyclic distance is the least absolute residue
of a difference modulo $M$; its radius-$R$ balls have at most
$2\lceil R\rceil+1$ sites. Thus the finite-range metric theorem supplies
$c_1,\eta>0$ and $C\ge0$ before the physical chain length and Hamiltonian
are chosen. With an original interaction bound $J_0$, take
$J=\max\{J_0,J_0C\}$ in the formulas above and define
\begin{align}
 M_i=\mathcal G_{c_1\ell\Delta,\Delta}
   \left(\sum_{j:\,d(i,j)\le\ell-R}H_j,\ H_i\right).
 \label{eq:hastings05_actual_periodic_terms}
\end{align}
Here $\mathcal G_{q,\Delta}(H',W)$ denotes the normalized Gaussian filter
of $W$ under $H'$. These are the actual truncated operators of (mid).
Their support, Hermiticity, norm and full-filter approximation estimates
are proved, not additional premises of the physical hierarchy theorem.
At $R=0$, the local terms commute and both filters equal $H_i$ exactly.
The uniform physical theorem therefore includes zero range and arbitrary
physical length.\footnote{Declarations:
\leanid{TNLean.MPS.AreaLaw.LocalSolution.periodicPhysicalMetricSpace},
\leanid{TNLean.MPS.AreaLaw.LocalSolution.card_periodicPhysical_metricRangeBall_le},
\leanid{TNLean.MPS.AreaLaw.LocalSolution.physicalPeriodicProjectionRadius},
\leanid{TNLean.MPS.AreaLaw.LocalSolution.periodicPhysicalTruncatedTerm},
\leanid{TNLean.MPS.AreaLaw.LocalSolution.exists_periodicPhysicalTruncatedTerm_localization}, and
\leanid{TNLean.MPS.AreaLaw.LocalSolution.exists_periodicPhysicalHierarchy}.}

The final physical theorem retains equation (low) as its local residual
premise, as in Section II of the source. It does not derive equation (low) for
an arbitrary vector in a degenerate ground space. Nor does it assert
that the additional closing window preserves a prescribed conventional
matrix product bond dimension; the source's representation is the tree
of nested retained subspaces.
\bibliographystyle{unsrt}
\bibliography{references}
\end{document}
